\chapter{Introduction}
\label{ch:introduction}

%% ── 1.1 ─────────────────────────────────────────────────────────────────────
\section{Context: LoRaWAN for Indoor Environmental Monitoring}
\label{sec:context}

% SET THE SCENE: LoRaWAN is everywhere for indoor sensing.
% Indoor propagation is complex — walls, people, air conditions.
% Accurate path loss modeling is critical for network planning and resource optimization.
% Transition: "A prior measurement campaign at [building] established that
%   environmental factors significantly affect LoRaWAN path loss indoors."

%% ── 1.2 ─────────────────────────────────────────────────────────────────────
\section{The Prior Project: Centralized Path Loss Modeling}
\label{sec:prior-project}

% THIS IS THE CONTINUATION ANCHOR — be very specific.
% Describe the prior project in detail:
%   - 6 MKR WAN 1310 nodes with BME280 + SCD4x + SPS30 sensors
%   - Deployed across indoor office (stairwell, corridor, kitchen, etc.)
%   - 8-month campaign: September 2024 — May 2025
%   - Collected 1,715,869 validated rows (aggregated from 2.3M raw)
%   - Features: pressure, CO2, temperature, humidity, PM2.5
%   - Radio metrics: RSSI, SNR, spreading factor, expected path loss (exp_pl)
%   - Trained centralized ML models: MLR, Random Forest, XGBoost
%   - Best result: XGBoost R² ≈ 0.93 for path loss regression
%   - Knowledge distillation: XGBoost → lightweight neural network

% INCLUDE TABLE: Prior project centralized ML results
\begin{table}[H]
  \centering
  \begin{tabular}{lcc}
    \toprule
    \textbf{Model} & \textbf{$R^2$} & \textbf{RMSE (dB)} \\
    \midrule
    Multiple Linear Regression & [XX] & [XX] \\
    Random Forest              & [XX] & [XX] \\
    XGBoost                    & 0.93 & [XX] \\
    \bottomrule
  \end{tabular}
  \caption{Centralized path loss regression results from the prior project.}
  \label{tab:prior-results}
\end{table}

% STATE THE LIMITATION:
% "Despite strong predictive performance, the centralized architecture suffered
%  from three fundamental limitations that motivate this thesis."

%% ── 1.3 ─────────────────────────────────────────────────────────────────────
\section{Problem Statement}
\label{sec:problem-statement}

% THREE CONCRETE PROBLEMS with the prior centralized system:

% 1. BANDWIDTH: each device transmitted 18 bytes of raw sensor data every 60 seconds,
%    producing 1440 messages/day × 6 nodes = 8640 msgs/day = ~156 KB/day.
%    Under EU868 1% duty cycle, this approaches regulatory limits.

% 2. PRIVACY: raw environmental data (temperature, CO2, humidity) leaves the device.
%    In occupied buildings, these readings can reveal occupancy patterns, activity,
%    and working habits — a potential GDPR concern.

% 3. ADAPTABILITY: the model was trained once on historical data and deployed statically.
%    It cannot adapt to seasonal changes, furniture rearrangement, building modifications,
%    or occupancy shifts that alter indoor propagation characteristics over time.

% CONCLUDE: "These limitations motivate a paradigm shift from centralized to federated
%   edge intelligence."

%% ── 1.4 ─────────────────────────────────────────────────────────────────────
\section{Proposed Solution: Federated TinyML}
\label{sec:proposed-solution}

% DESCRIBE THE NEW APPROACH in one paragraph:
% - Each device runs a compact neural network (~57 parameters, <500 bytes)
% - Predicts expected path loss from local environmental observations
% - Trains locally on buffered samples using simplified on-device SGD
% - Sends only quantized weight updates (~52 bytes) per FL round via LoRaWAN
% - Server aggregates using FedAvg and redistributes an improved global model

% CONTRAST WITH TORRES SANCHEZ et al.:
% "Torres Sanchez et al.~\cite{torressanchez2024federated} recently demonstrated
%  the feasibility of FL over LoRaWAN for IIoT anomaly detection, achieving
%  an F1 score of 94.77\% using virtual clients in a Flower-based simulation.
%  Our work extends this direction to environment-driven path loss regression
%  on physically deployed sensor nodes, using real measurement data and
%  deploying TFLite Micro models on actual MKR WAN 1310 hardware."

%% ── 1.5 ─────────────────────────────────────────────────────────────────────
\section{Research Objectives}
\label{sec:objectives}

The primary objectives of this thesis are:
\begin{enumerate}
  \item To design and implement a \gls{fl} framework for path loss regression across
        multiple \gls{lorawan} edge devices without sharing raw sensor data.
  \item To develop a \gls{tinyml} regression model compact enough to perform both
        inference and local training on the Arduino MKR~WAN~1310
        (SAMD21 Cortex-M0+, \SI{32}{\kilo\byte} \gls{sram}, \SI{256}{\kilo\byte} Flash).
  \item To evaluate the accuracy trade-off between centralized and federated path loss
        regression using the real-world measurement dataset, quantified by $R^2$ and
        \gls{rmse}.
  \item To quantify the communication efficiency of the federated approach compared to
        raw data transmission under EU868 duty-cycle constraints.
  \item To validate the full system pipeline on real hardware as a proof-of-concept
        prototype, directly extending the prior project's deployment.
\end{enumerate}

%% ── 1.6 ─────────────────────────────────────────────────────────────────────
\section{Research Questions}
\label{sec:research-questions}

\begin{description}
  \item[\textbf{RQ1}] Can a federated \gls{tinyml} model achieve comparable path loss
    prediction accuracy ($R^2$) to a centralized model trained on the full dataset?

  \item[\textbf{RQ2}] How does non-\gls{iid} data distribution---arising from devices
    placed in physically distinct rooms---affect federated model convergence?

  \item[\textbf{RQ3}] What bandwidth reduction does transmitting model weight updates
    ($\approx$\,\SI{52}{\byte}/round) achieve over raw sensor data transmission
    (\SI{18}{\byte}$\times$1440/day)?

  \item[\textbf{RQ4}] Can the complete \gls{fl} pipeline (inference, local training,
    weight quantization, LoRaWAN communication) operate within the memory and compute
    constraints of the MKR~WAN~1310?
\end{description}

%% ── 1.7 ─────────────────────────────────────────────────────────────────────
\section{Contributions}
\label{sec:contributions}

This thesis makes the following contributions:
\begin{enumerate}
  \item \textbf{First \gls{fl} framework for environment-driven path loss regression
    on \gls{lorawan}.} No prior work combines federated learning with TinyML for path
    loss prediction on resource-constrained LoRaWAN devices. Torres Sanchez et
    al.~\cite{torressanchez2024federated} demonstrated FL feasibility over LoRaWAN for anomaly
    detection; we extend this to continuous regression on real environmental data.

  \item \textbf{Evaluation on real-world measurements.} Unlike simulation-only studies,
    our federated evaluation uses 1,715,869 measurements from six physically deployed
    MKR~WAN~1310 nodes, with non-\gls{iid} data partitioning by actual device location.

  \item \textbf{Direct comparison across three paradigms.} We present a three-way
    comparison: (1)~the prior project's centralized XGBoost ($R^2 \approx 0.93$),
    (2)~a centralized neural network baseline, and (3)~the federated neural network.

  \item \textbf{Working hardware prototype.} A \gls{tflite} Micro model deployed on a
    real MKR~WAN~1310 demonstrates that the pipeline fits within \SI{32}{\kilo\byte}
    SRAM, complementing the simulation-based evaluation.
\end{enumerate}

%% ── 1.8 ─────────────────────────────────────────────────────────────────────
\section{Thesis Outline}
\label{sec:outline}

The remainder of this thesis is organized as follows.
Chapter~\ref{ch:background} presents the theoretical foundations of LoRaWAN, path loss
modeling, federated learning, and TinyML.
Chapter~\ref{ch:related-work} reviews related work and identifies the research gap.
Chapter~\ref{ch:methodology} presents the full methodology: system design, dataset,
model architecture, training pipeline, federated learning protocol, simulation setup,
edge firmware, and aggregation server.
Chapter~\ref{ch:results-discussion} presents the experimental results and discussion.
Chapter~\ref{ch:conclusion} concludes the thesis and outlines future work.
